\begin{frame}{학습 곡선: ``방향''을 빠르게 잡아 버린다}
  \begin{columns}[T]
    \begin{column}{0.5\textwidth}
      \centering
      \includegraphics[width=\linewidth]{ch12_2_reward_loss_curve.png}
      {\footnotesize x=에폭, y=에폭당 평균 손실 $-\log\sigma(r_w-r_l)$.
      처음 수십 에폭 급락, 이후 거의 수평.}
    \end{column}
    \begin{column}{0.48\textwidth}
      \begin{itemize}
        \item 300 쌍, 200 에폭: 처음 \textbf{급격히} 떨어지다가
          어느 시점부터 \textbf{수평} — ``부드러움이 더 가치
          있다''는 \textbf{방향}을 빠르게 잡은 것.
        \item 곡선이 0까지 안 내려가는 이유: \textbf{가까운} 궤적
          쌍은 진짜 보상에서도 차이가 작아 $\sigma(r_w-r_l)$이 1에
          완전히 도달 못 함.
        \item 손실의 \textbf{절대값}보다 \textbf{하향 추세}를 읽으라.
      \end{itemize}
    \end{column}
  \end{columns}
\end{frame}
